\documentclass[11pt,a4paper]{article}
\usepackage[T1]{fontenc}
\usepackage[utf8]{inputenc}
\usepackage{amsmath,amssymb,amsthm}
\usepackage[margin=1in]{geometry}
\usepackage[colorlinks=true,linkcolor=blue,urlcolor=blue]{hyperref}
\theoremstyle{plain}
\newtheorem{theorem}{Theorem}
\newtheorem{lemma}[theorem]{Lemma}
\newtheorem{proposition}[theorem]{Proposition}
\newtheorem{corollary}[theorem]{Corollary}
\theoremstyle{definition}
\newtheorem{remark}[theorem]{Remark}

\title{Recovering the unwritten recurrences of the table OEIS A236746}
\author{Adrian Perez Fontelles\\ \small Independent researcher}
\date{2 September 2026}
\begin{document}
\maketitle

\begin{abstract}
OEIS A236746 is a two-dimensional table $T(n,k)$ whose empirical recurrences are, for several of
its lines, given only as an ORDER --- ``k=2: [order 18]'' and the like --- with the
coefficients in a linked file rather than in the entry. Those conjectures can still be settled,
and the recurrences recovered. A line of the table is a fixed-width or fixed-height array
count, hence a walk count on finitely many vertices, hence satisfies some monic recurrence of
order at most that number; Berlekamp--Massey on twice as many exact terms returns the MINIMAL
one. Where its order equals the order the entry states --- and where the same holds after the
first few terms are dropped --- any recurrence of that order the line satisfies has a
characteristic polynomial that is a multiple of the minimal one and of equal degree, hence is
that polynomial. This note settles column 2, column 3, column 4.
\end{abstract}

\noindent\small 2020 Mathematics Subject Classification. 05A15, 11B37, 15A18.\normalsize

\section{The table and the unwritten conjectures}

OEIS A236746 is ``T(n,k)=Number of (n+1)X(k+1) 0..2 arrays with the upper median plus the lower median of every 2X2 subblock equal.''. It is read by antidiagonals and begins
\[
81, 267, 267, 833, 1411, 833, 2907, 7377,\ \dots
\]
The entry carries, among its empirical claims:
\begin{quote}\small
k=2: [order 18]\\
k=3: [order 38]\\
k=4: [order 86]
\end{quote}
As of the ``Last modified'' line on the live entry (Jul 23 2025 09:31:42, revision 6) these are still
recorded as empirical, and the recurrences themselves are not in the entry text.

\section{Why a line of the table is a walk count}

Fix the index that the line fixes. The entry defines $T(n,k)$ as a number of arrays of a shape
determined by $n$ and $k$, subject to a condition local to a bounded window of consecutive
lines. Substituting the fixed index into the entry's own wording turns the definition into an
ordinary one-parameter array count, and for such a count the admissible windows are the
vertices of a finite digraph and an array is a walk. So the line is a walk count on $S$
vertices, and by the Cayley--Hamilton theorem it satisfies a monic integer recurrence of order
at most $S$.

\section{Recovering the recurrences}

\begin{lemma}\label{lem:bm}
A sequence known to satisfy some linear recurrence of order at most $S$ is determined, with its
minimal recurrence, by its first $2S$ terms; Berlekamp--Massey applied to them returns that
minimal recurrence exactly.
\end{lemma}

\subsection*{Column $k=2$}
Substituting the fixed index into the entry's wording gives an ordinary one-parameter array
count, a walk on $S=93$ vertices. Berlekamp--Massey on $2S=186$ exact terms
returns a minimal recurrence of order $18$ --- the order the entry states --- with
integer coefficients
\begin{center}\small\begin{tabular}{cccccccc}
$c_{1}$ & $17$ & $c_{6}$ & $-14868$ & $c_{11}$ & $473554$ & $c_{16}$ & $-171072$ \\
$c_{2}$ & $-78$ & $c_{7}$ & $37484$ & $c_{12}$ & $-502900$ & $c_{17}$ & $157120$ \\
$c_{3}$ & $-142$ & $c_{8}$ & $39431$ & $c_{13}$ & $-310088$ & $c_{18}$ & $-32640$ \\
$c_{4}$ & $1902$ & $c_{9}$ & $-204023$ & $c_{14}$ & $766120$ &  &  \\
$c_{5}$ & $-2258$ & $c_{10}$ & $60026$ & $c_{15}$ & $-297584$ &  &  \\
\end{tabular}\end{center}
so that $a(m)=\sum_{i=1}^{18}c_i\,a(m-i)$ for every $m$. Removing the first
$18+4$ terms and repeating gives order $18$ again, which is what the
identification below needs.

\subsection*{Column $k=3$}
Substituting the fixed index into the entry's wording gives an ordinary one-parameter array
count, a walk on $S=243$ vertices. Berlekamp--Massey on $2S=486$ exact terms
returns a minimal recurrence of order $38$ --- the order the entry states --- with
integer coefficients
\[
(c_1,\dots,c_{38})=(29,\ -236,\ -530,\ 16058,\ -44099,\ -359428,\ 2037666,\ 2641959,\ -38207537,\ 26566108,\ 394980672,\ -749280468,\ -2366139159,\ 7499380996,\ 7313304008,\ -43398468452,\ -1170978370,\ 160201297848,\ -83692886880,\ -390363727264,\ 354528700808,\ 637126357152,\ -786883396000,\ -700593013632,\ 1098115601088,\ 520503186176,\ -1014936927488,\ -264788255744,\ 628983984128,\ 97800536064,\ -257790742528,\ -29839360000,\ 66734374912,\ 7884767232,\ -9802874880,\ -1429733376,\ 613416960,\ 113246208).
\]
so that $a(m)=\sum_{i=1}^{38}c_i\,a(m-i)$ for every $m$. Removing the first
$38+4$ terms and repeating gives order $38$ again, which is what the
identification below needs.

\subsection*{Column $k=4$}
Substituting the fixed index into the entry's wording gives an ordinary one-parameter array
count, a walk on $S=657$ vertices. Berlekamp--Massey on $2S=1314$ exact terms
returns a minimal recurrence of order $86$ --- the order the entry states --- with
integer coefficients
\[
(c_1,\dots,c_{86})=(56,\ -971,\ -152,\ 190873,\ -1661887,\ -10009206,\ 213828146,\ -312175856,\ -12610205483,\ 66713398446,\ 384188036780,\ -4069754544077,\ -2970246585544,\ 145450214494923,\ -260085066899947,\ -3389529900369997,\ 13437100535257914,\ 49901573197783878,\ -364155800303484757,\ -309048914527576159,\ 6693706416724637849,\ -5293405206431529293,\ -88446472645102617151,\ 189703042852733857249,\ 832480918929966534951,\ -3147931526385388195938,\ -4982867947477027772269,\ 35744713444977948535612,\ 6602349093447677536009,\ -301209810911293867384168,\ 238699555746690517592636,\ 1917867056603366392212474,\ -3250879857566599449169865,\ -9021390986329552113938280,\ 25702420173974928592296374,\ 28130870271634328062490137,\ -147062988088131091098368746,\ -26482957345743932638765509,\ 645125392261628761703284995,\ -304952979906921662062790849,\ -2202780261306128925726760480,\ 2329867427067214688914531910,\ 5787794445026478112504831387,\ -10041438664471031213750311635,\ -11078842104635443948380076365,\ 31468475919075686426485822419,\ 12354722441315253154852427645,\ -76519317526070181580388021284,\ 6133192123440808923868375827,\ 147531434786839346425260796567,\ -67607820835288691508451652323,\ -225963172752204105161432620551,\ 185450411753089956307334138380,\ 270558574911593990653527754082,\ -338725158984611952959386745108,\ -241102032962494027581075625768,\ 467511584304256743355102024856,\ 135477959681789037882763204724,\ -507016704763296591995520520240,\ -1807913563854425524833517816,\ 437526343153511336076569994464,\ -95174312359171190384206308032,\ -300348193561033103745352744768,\ 123182847888821401799137972864,\ 162396678032868024066180336384,\ -97810180468021909967610445824,\ -67701854391562253369639063552,\ 56163127676633147343642025984,\ 20872466437286802216809275392,\ -24342881261878518215026139136,\ -4318515969671209454095187968,\ 8050283096077472683816943616,\ 403752404210830088816361472,\ -2025583667668892646967345152,\ 70664439126020602398769152,\ 383578725086594313962913792,\ -34184188020431730845417472,\ -53678080702240805156290560,\ 6310719444052906113761280,\ 5392813591250939456520192,\ -657097221731299268493312,\ -369941503469134156922880,\ 37829871798266458275840,\ 15652394185300063027200,\ -940704535119554150400,\ -311136828310683648000).
\]
so that $a(m)=\sum_{i=1}^{86}c_i\,a(m-i)$ for every $m$. Removing the first
$86+4$ terms and repeating gives order $86$ again, which is what the
identification below needs.

\begin{theorem}
For each line above, the displayed recurrence holds for every index, and it is the recurrence
the entry records.
\end{theorem}

\begin{proof}
It holds by Lemma~\ref{lem:bm}. For the identification, the entry asserts that the line
satisfies SOME recurrence of the stated order, possibly only from some index on. Let $q$ be its
characteristic polynomial and $q_{\min}$ the minimal polynomial of the tail. Then
$q_{\min}\mid q$, and the second Berlekamp--Massey run --- on the line with its first few
terms removed --- shows $\deg q_{\min}$ equals the stated order, which is $\deg q$. Both are
monic, so $q=q_{\min}$.
\end{proof}

\section{Verification}

Three checks.

First, each line's model was compared against the entry's own published values for that line,
taken from the antidiagonal DATA read in either orientation and from the ``Table starts''
block, and agrees at every available term. This is the only point at which reading the entry's
English enters, and a misreading gives different counts.

Second, each recovered recurrence was evaluated directly on the model's values, with no
Berlekamp--Massey involved, and holds at every index tested.

Third, each order was computed twice by independent routes --- modulo a $61$-bit prime, where
every intermediate is a machine word, and again in exact rational arithmetic --- and once more
on the line with its first few terms dropped, which is what the identification requires. All
agree.

\begin{thebibliography}{9}
\bibitem{oeis} The OEIS Foundation, \emph{The On-Line Encyclopedia of Integer
Sequences}, \url{https://oeis.org}, 2026. Sequence A236746.
\bibitem{massey} J.~L.~Massey, Shift-register synthesis and BCH decoding, \emph{IEEE
Trans. Inform. Theory} \textbf{15} (1969), 122--127.
\bibitem{stanley} R.~P.~Stanley, \emph{Enumerative Combinatorics}, Volume 1, 2nd ed.,
Cambridge University Press, 2011. (Transfer-matrix method, Section 4.7.)
\bibitem{horn} R.~A.~Horn and C.~R.~Johnson, \emph{Matrix Analysis}, 2nd ed., Cambridge
University Press, 2013.
\end{thebibliography}

\end{document}
